% Source ownership:
% - Oliwier: energy models, OPLS-AA, topology, delta_energy
% - Manel: Geweke runtime equilibration and MC workflow
% - Arthur: integration / prose cleanup

% - Keep modular because Oliwier will still update references

\section{Energy Evaluation and Monte Carlo Engine}
\label{sec:energy_mc}

\subsection{Energy Calculations}

The energy evaluation modules, \texttt{energy.f90} and \texttt{energy\_all\_atoms.f90}, were designed and implemented to compute the total potential energy of the chain and the $\Delta E$ increments required by the Metropolis acceptance/rejection step.

\subsubsection{United-Atom (UA) Model \& TraPPE-UA}

\textbf{Trigonometric bypass in dihedral calculations.}

The standard approach to evaluating a dihedral angle requires an
\texttt{acos} or \texttt{atan2} call, both of which are transcendental
functions with significant CPU cost relative to basic arithmetic.
Instead, \texttt{compute\_cos\_dihedral} computes $\cos\phi$ directly
from the dot product of the normal vectors of the two planes defined by
the carbon backbone:

\begin{equation}
  \cos\phi = \frac{\mathbf{n}_1 \cdot \mathbf{n}_2}{|\mathbf{n}_1|\,|\mathbf{n}_2|},
  \qquad
  \mathbf{n}_1 = \mathbf{b}_1 \times \mathbf{b}_2,\quad
  \mathbf{n}_2 = \mathbf{b}_2 \times \mathbf{b}_3
\end{equation}

where $\mathbf{b}_i$ are successive bond vectors. A guard clause handles
collinear atoms (degenerate normals) by setting $\cos\phi = 1$ as a
fallback. The function is declared \texttt{pure}, signalling to the compiler
that it has no side effects, which allows it to be called safely
from within parallel regions and \texttt{forall} constructs.

\textbf{Polynomial expansion for TraPPE-UA.}

The TraPPE-UA torsional potential~[3] is:

\begin{equation}
  U(\phi) = c_1(1+\cos\phi) + c_2(1-\cos 2\phi) + c_3(1+\cos 3\phi)
\end{equation}

Substituting $y = \cos\phi$ and applying the identities
$\cos 2\phi = 2y^2-1$ and $\cos 3\phi = 4y^3-3y$ gives a polynomial
in $y$:

\begin{equation}
  U(y) = c_1(1+y) + c_2(2-2y^2) + c_3(1-3y+4y^3)
\end{equation}

Since $y = \cos\phi$ is already available from the cross-product step,
the torsional energy is evaluated with only multiplications and
additions -- no further transcendental calls are needed in the MC loop.

\subsubsection{Explicit Hydrogens and the OPLS-AA Force Field}

The all-atom model requires handling a substantially more complex
topology. The OPLS-AA force field~[4] was adopted, with atom-specific
Lennard-Jones parameters for carbon and hydrogen
($\sigma_\mathrm{CC}=3.50\,\text{\AA}$,
$\varepsilon_\mathrm{CC}=0.066\,\text{kcal/mol}$;
$\sigma_\mathrm{HH}=2.50\,\text{\AA}$,
$\varepsilon_\mathrm{HH}=0.030\,\text{kcal/mol}$; cross-terms from
geometric mixing rules).

\textbf{Effective backbone torsional potential.}

In OPLS-AA, rotating a single C--C bond changes nine distinct dihedral
angles: one C--C--C--C, four C--C--C--H, and four H--C--C--H. Computing
all nine at every MC step is expensive for a code designed to run
millions of steps. An \emph{effective backbone potential} was therefore
constructed by summing the OPLS-AA Fourier coefficients for all nine
dihedral types analytically, producing a single set of effective
coefficients that reproduce the combined torsional energy from the
C--C--C--C backbone coordinates alone:

\begin{equation}
  c_1^\mathrm{eff} = 0.870,\quad
  c_2^\mathrm{eff} = -0.0785,\quad
  c_3^\mathrm{eff} = 1.5075 \quad \text{(kcal/mol)}
\end{equation}

The OPLS-AA Lennard-Jones parameters are taken from the framework of
Sæther et al.~[4]; the effective coefficient derivation is an original
contribution of this work. The same polynomial form used for the UA
model is then applied with these coefficients, so the all-atom
torsional evaluation has the same arithmetic cost as the united-atom one.

The corresponding \texttt{torsion\_single} function in
\texttt{energy\_all\_atoms.f90} is:

\begin{lstlisting}
pure function torsion_single(cos_phi) result(e)
  double precision, intent(in) :: cos_phi
  double precision :: e, y, y2, y3

  y  = cos_phi
  y2 = y * y
  y3 = y2 * y

  e = opls_c1 * (1.0d0 + y)            &
    + opls_c2 * (2.0d0 - 2.0d0 * y2)   &
    + opls_c3 * (1.0d0 - 3.0d0 * y + 4.0d0 * y3)
end function torsion_single
\end{lstlisting}

\textbf{Topology mapping and exclusion matrix.}

Non-bonded Lennard-Jones interactions must be suppressed for atom pairs
separated by fewer than four bonds (1-2, 1-3, and 1-4), since those
interactions are already captured by the bonded energy terms. The
initialisation routine \texttt{init\_energy\_topology} assigns each
hydrogen to its parent carbon by geometric distance (C--H threshold
$1.2\,\text{\AA}$), records the backbone position index for every atom,
and builds a global Boolean matrix \texttt{is\_excluded}. Each pair is
then looked up in $O(1)$ per call, avoiding any bond traversal during
the MC loop.

The exclusion criterion uses a uniform backbone-position threshold: pair
$(i,j)$ is excluded when
$|\text{backbone\_pos}(i) - \text{backbone\_pos}(j)| < 4$.
For C--C pairs this is exact. For C--H and H--H pairs the actual bond
count exceeds the backbone-position difference by one and two
respectively, so the threshold over-excludes some 1-5 and 1-6
interactions. The torsional potential partially compensates, since its
parameterisation implicitly encodes the conformational barriers arising
from steric repulsion. The single-chain, vacuum context of the
simulation further limits the practical impact relative to a condensed
phase. Correcting the exclusion criterion to a bond-count-aware rule is
left for future work.

\textbf{Dynamic atom lists for $\Delta E$.}

At each pivot move only atoms downstream of the chosen bond change
position. Recomputing the full $O(N^2)$ LJ sum is therefore unnecessary.
Instead, the atoms are partitioned into \texttt{fixed\_list} and
\texttt{moved\_list} by comparing old and new coordinate arrays:

\begin{lstlisting}
n_fixed = 0
n_moved = 0
do i = 1, n_atoms
  if (sum((coords_new(i,:) - coords_old(i,:))**2) > 1.0d-12) then
    n_moved = n_moved + 1
    moved_list(n_moved) = i
  else
    n_fixed = n_fixed + 1
    fixed_list(n_fixed) = i
  end if
end do
\end{lstlisting}

Only cross-interactions between fixed and moved atoms are re-evaluated,
reducing the $\Delta E$ cost to $O(N_\mathrm{fixed} \times
N_\mathrm{moved})$. A secondary \texttt{pair\_type\_matrix} routes each
valid pair to the pre-computed C--C, C--H, or H--H LJ parameters in a
single array lookup.

\subsection{Pivot Move and Monte Carlo Step}
 
The conformational sampling mechanism used throughout the project is the
pivot move. Rather than perturbing the whole chain diffusely, the
algorithm selects one internal C--C bond and rotates the entire tail of
the polymer as a rigid body around that bond axis. This generates large
collective rearrangements and allows the chain to explore configuration
space much more efficiently than local coordinate displacements.

The rigid-body rotation is implemented with Rodrigues' rotation formula.
If \(\mathbf{v}\) is the position of an atom relative to the pivot
point, \(\mathbf{k}\) is the unit vector along the selected bond, and
\(\theta\) is the random rotation angle, the rotated vector is
\[
\mathbf{v}_{\mathrm{rot}}
=
\mathbf{v}\cos\theta
+
(\mathbf{k}\times\mathbf{v})\sin\theta
+
\mathbf{k}(\mathbf{k}\cdot\mathbf{v})(1-\cos\theta).
\]
This construction preserves bond lengths and bond angles while applying
a clean rigid-body rotation to the selected downstream segment.

\subsubsection{Handling $sp^3$ Hybridisation in All-Atom Mode}

In all-atom mode, the same geometric transformation must be applied not
only to the carbon backbone but also to the hydrogens attached to every
rotating carbon. Because these hydrogens represent an approximately
tetrahedral $\mathrm{sp}^3$ environment, any inconsistency between the
carbon motion and the attached hydrogen motion would immediately distort
the local geometry and generate unphysical configurations. The AA
rotation routine therefore identifies the hydrogens attached to each
moving carbon and applies the exact same pivot point, axis, and angle to
them, using a C--H bond threshold of $1.2\,\text{\AA}$:
 
\begin{lstlisting}
if (explicit_h) then
  ! Hydrogens are stored after the carbon backbone in the coords array
  do i = n_carbons + 1, n_atoms
    do j = k + 1, n_carbons
      ! Identify if hydrogen i is bonded to the moving carbon j
      v = coords(i, :) - coords(j, :)
      if (vnorm(v) < 1.2d0) then
        ! Apply Rodrigues' rotation relative to the pivot point
        v = coords(i, :) - pivot
        dot_uv = sum(axis * v)
        v_rot = v * cos_p + cross(axis, v) * sin_p &
              + axis * dot_uv * (1.0d0 - cos_p)
        coords_new(i, :) = pivot + v_rot
        exit  ! Hydrogen found, move to the next atom
      end if
    end do
  end do
end if
\end{lstlisting}

\subsubsection{The MC Step and Metropolis Criterion}

Each call to \texttt{mc\_step} performs one trial move. First, a random
internal bond index $k$ is drawn uniformly from $[1,\, N_C - 2]$, and a
random rotation angle $\phi$ is sampled uniformly from
$[-\Delta\phi_{\max},\, \Delta\phi_{\max}]$; the parameter
\texttt{max\_delta} is tuned to achieve a reasonable acceptance rate.
Second, a trial configuration is generated by rotating the corresponding
tail segment while leaving the current coordinates unchanged. Third, the
change in energy is computed using the model-appropriate $\Delta E$
routine: \texttt{delta\_energy\_aa} for all-atom mode and
\texttt{delta\_energy\_ua} for united-atom mode, dispatched via the
\texttt{explicit\_h} flag. Fourth, the trial move is accepted or
rejected using the Metropolis criterion.

The Metropolis rule is standard. If the move lowers the total energy, it
is accepted unconditionally. Otherwise, it is accepted with probability
\[
P_{\mathrm{acc}} = \exp(-\beta \Delta E),
\qquad
\beta = \frac{1}{k_B T}.
\]
If accepted, the trial coordinates replace the current coordinates and
the total, Lennard-Jones, and torsional energies are updated in place.
If rejected, the trial geometry is discarded and the original state is
retained. Repeating this procedure generates a Markov chain whose
stationary distribution is the desired Boltzmann distribution.
 
\begin{lstlisting}
! d. Metropolis acceptance criterion
call random_number(random_value)
if (dE < 0.0d0 .or. random_value < exp(-beta * dE)) then
  coords  = coords_new
  E_total = E_total + dE
  E_lj    = E_lj    + dE_lj
  E_tors  = E_tors  + dE_tors
  accepted = .true.
end if
\end{lstlisting}
 
\subsection{Geweke Convergence Detection}

Detecting equilibration in a Markov chain is non-trivial because
successive configurations are strongly autocorrelated. Earlier versions
of the workflow relied on a fixed burn-in length and a posterior
post-processing check based on a Fast-Fourier Transform (FFT) script to
estimate the integrated autocorrelation time $\tau_\mathrm{int}$, but
this was often too rigid and could either discard too much data or fail
to detect persistent transients, sometimes reporting that the production
portion was 0\% of the 10M steps. To make equilibration detection
adaptive, the serial Monte Carlo engine was extended with an
in-simulation Geweke diagnostic~[1].

The Geweke test compares the mean of an early window of the trajectory
with the mean of a late window of the same trajectory. In our
implementation, the monitored quantity is the total energy. If the chain
has equilibrated, both windows should be consistent with the same
stationary distribution. The resulting $z$-score is
\[
z = \frac{\bar{x}_A - \bar{x}_B}
        {\sqrt{\mathrm{SE}_A^2 + \mathrm{SE}_B^2}},
\]
where $\bar{x}_A$ and $\bar{x}_B$ are the mean energies in the early
and late windows, and $\mathrm{SE}_A$ and $\mathrm{SE}_B$ are their
standard errors.

A crucial point is that the standard errors cannot be estimated from the
raw sample variance as if the measurements were independent. Monte Carlo
trajectories are correlated, so such a treatment would underestimate the
uncertainty and produce false positives. To account for this, the code
uses batch means. The window is partitioned into contiguous blocks, the
mean of each block is computed, and the variance of those block means is
used as an estimator of the uncertainty in the presence of
autocorrelation:

\begin{equation}
  SE_A^2 = \frac{1}{b_A(b_A-1)}\sum_{i=1}^{b_A}\left(\bar{x}_i - \bar{x}_A\right)^2
\end{equation}

where $\bar{x}_i$ is the mean of block $i$ and $\bar{x}_A$ is the mean
of the full window. $SE_B^2$ is computed analogously over window B.

The implementation uses an accumulated buffer of $N=300$ energy samples.
The early window contains the first $f_A=10\%$ of the buffer ($n_A=30$
samples, $b_A=\lfloor\sqrt{30}\rfloor=5$ blocks) and the late window
contains the last $f_B=50\%$ ($n_B=150$ samples,
$b_B=\lfloor\sqrt{150}\rfloor=12$ blocks). Because block means of
sufficiently large blocks behave approximately as independent, this
estimator correctly absorbs the autocorrelation structure of the chain.

Under the null hypothesis of stationarity, $z$ follows a standard
normal distribution. Equilibrium is declared when
$|z| < z_\mathrm{crit} = 1.96$ (95\% confidence level) on
$n_\mathrm{consec} = 3$ consecutive evaluations, preventing false
positives caused by temporary plateaus. The test is re-evaluated every
\texttt{eval\_freq} synchronisation points once the buffer is full,
so the simulation can stop the equilibration phase as soon as a
statistically stable regime has been reached. The production phase then
natively restarts step counting from 1.

